\documentclass[11pt]{article}

\usepackage[T1]{fontenc}
\usepackage{amsmath,amssymb}
\usepackage{booktabs}
\usepackage{geometry}
\usepackage{listings}
\usepackage{xcolor}
\usepackage{hyperref}
\geometry{margin=1in}

\lstdefinelanguage{FWATS3}{
  morekeywords={val,lam,fix,if,then,else,let,in,end,dec,imp},
  sensitive=true,
  morecomment=[l]{\#},
  morestring=[b]"
}
\lstset{language=FWATS3,basicstyle=\ttfamily\small,columns=fullflexible,
  keepspaces=true,frame=single,breaklines=true}

\newcommand{\FWATS}{\textsc{Fwats3}}
\newcommand{\Var}{\mathit{Var}}
\newcommand{\Con}{\mathit{Con}}
\newcommand{\Imp}{\mathit{Imp}}
\newcommand{\matchsto}{\mathrel{\Downarrow_{\!m}}}

\title{FWATS3: A Typed Core Language with Explicit Template Dispatch}
\author{Draft}
\date{\today}

\begin{document}
\maketitle

\begin{abstract}
\FWATS{} is a small expression language in which ordinary values and
parameterized constants occupy distinct namespaces. A declaration gives a
constant a type scheme, while one or more separately written implementations
associate static patterns with expressions. A use supplies its static
arguments explicitly; the implementation is selected by directional pattern
matching and its body is instantiated before evaluation. The language combines
this mechanism with annotated functions, recursion, tuples, sequential declarations,
and a structural type checker. This paper introduces the language's central
distinction, gives its core declaration and typing model, and describes the
scope of the current interpreter. The present implementation is an executable
prototype: implementation availability is a run-time property, dispatch is
ordered, and the checker does not establish termination or implementation
coverage.
\end{abstract}

\section{Introduction}

Parameterized computation is often expressed with polymorphic functions,
macros, type classes, or compiler specialization. These mechanisms differ in
when a parameter is interpreted and how a concrete implementation is selected.
\FWATS{} isolates one particular arrangement: a programmer declares the type
of a named constant, writes implementations separately from that declaration,
and explicitly supplies static arguments at each use. Implementations are
selected by matching those arguments against ordered static patterns. The
selected expression is then evaluated as an ordinary expression in a
lexically captured value environment.

This arrangement makes two forms of name resolution visible in the language.
An unprefixed expression name denotes a dynamically bound variable; a name
prefixed by \texttt{@} denotes a template constant. The former is resolved in
the value environment, while the latter is resolved through a collection of
implementation closures. A template constant is therefore neither an ordinary
polymorphic value nor a macro that expands without regard to runtime scope.
Static arguments choose an implementation; the resulting value can then be
used as a conventional, monomorphic value, including as a function argument.

The main example is the identity template:
\begin{lstlisting}
dec identity[a]: a -> a;
imp{a} identity[a] = lam(x: a) => x;
@identity[sint](42)
\end{lstlisting}
The declaration introduces a signature parameterized by the static variable
\(a\). The implementation binds its own variable, also written \(a\), and
matches that variable against the use-site argument \texttt{sint}. The
implementation body is checked at type \texttt{sint -> sint} after the
signature is instantiated, and evaluation returns \texttt{42}. The parameter
names in the declaration and implementation are independently bound; their
spelling does not connect the two scopes.

\FWATS{} is designed to make the boundary between type-level selection and
ordinary evaluation explicit and inspectable. Its prototype includes a
recursive-descent parser, a structural type checker, and an interpreter over
the same syntax tree. The central contribution of this draft is a precise
account of that boundary. It is not a claim of a complete polymorphic type
system: the current checker uses explicit static arguments, supports no
first-class polymorphic template values, and checks each implementation
against its signature without proving that every use has a matching
implementation.

\paragraph{Contributions.}
This introduction presents (1) a source-level distinction between variables
and template constants, (2) a declaration discipline separating signatures
from pattern-indexed implementations, and (3) an execution model in which
implementation bodies capture value bindings while resolving nested template
uses in the active template environment. It also records the present
implementation's limits so that the language design is not confused with
properties the prototype does not guarantee.

\section{Language overview}

\subsection{Expressions and types}

The ordinary expression fragment contains integer, Boolean, and string
literals; annotated single-argument functions; recursive functions; function
application; conditionals; tuples and projections; sequential \texttt{let}
declarations; and primitive arithmetic and comparisons. An arrow type
\(T_1 \to T_2\) describes a function taking one value of type \(T_1\) and
returning a value of type \(T_2\). Multiple arguments are represented by one
tuple argument. Types include named constructors, functions, and tuples. Static
variables are introduced explicitly by declaration parameter lists; there is
no type-variable inference.

The core type constructors include \texttt{sint}, \texttt{bool},
\texttt{string}, and unit \texttt{()}, together with unary constructors such
as \texttt{list(T)}, \texttt{optn(T)}, \texttt{strm(T)}, and
\texttt{stcn(T)}. The latter constructors are nominal. The implementation
uses structural equality for types, so constructor names and ordered
arguments must agree.

\subsection{Signatures and implementations}

A signature has the form
\[
  \texttt{dec } c[p_1,\ldots,p_n]\texttt{: }T.
\]
It introduces a template constant \(c\) and binds the static parameters
\(p_i\) in \(T\). An implementation has the form
\[
  \texttt{imp}\{q_1,\ldots,q_m\}\ c[P_1,\ldots,P_n]
  \texttt{ = }e.
\]
Its implementation parameters \(q_j\) scope over the patterns \(P_i\) and
body \(e\). There is one pattern for each signature parameter, but the
implementation may bind a different number of variables. For example,
\texttt{identity[(x,y)]} can bind two implementation variables in one pattern
for a signature with one static parameter.

The checker first looks up the signature for \(c\), substitutes the
implementation patterns for its signature parameters, then checks \(e\)
against the resulting type. Implementation parameters are rigid type
variables during this check. Every such parameter must occur in a pattern;
the checker does not guess a concrete type to make an implementation body
check. Thus a generic implementation that uses integer addition directly
cannot type-check at an unconstrained type variable; it must use a separately
declared operation template or specialize its pattern.

\subsection{Explicit specialization and dispatch}

A use supplies static arguments in brackets:
\[
  \texttt{@}c[T_1,\ldots,T_n].
\]
The checker instantiates the declared signature with these arguments. It does
not select an implementation. During evaluation, the interpreter scans
registered implementations of \(c\) from newest to oldest. It matches the
entire pattern list against the entire argument list with one shared matching
context. A pattern variable binds to a type; repeated occurrences must match
structurally equal types. Constructors, function types, and tuples match
componentwise. The first successful match is selected. The body is copied,
substituted with the bindings, and evaluated.

This dispatch rule makes declaration order observable when patterns overlap.
There is no ranking by specificity, coherence check, or requirement that
implementations be exhaustive. Static arguments must be explicit and ground
at runtime. An ordinary function value produced by specialization can be
passed to another function, but it has a fixed type: an unspecialized template
is not a first-class polymorphic value.

\section{Static checking and execution}

\subsection{Structural type checking}

The checker synthesizes a type for each expression. Function parameters are
typed by their annotations; applications require the argument type to match
the function domain; conditionals require a Boolean condition and equal
branch types; tuples synthesize componentwise types; projections require an
in-bounds tuple index. Recursive functions include an annotated result type,
which the body must satisfy. Declarations are processed sequentially, so each
declaration may refer to preceding declarations.

The signature checker and the implementation dispatcher answer different
questions. Given a declared signature for \(c\), a static application can be
well-typed even when no implementation matches its arguments. Conversely,
each registered implementation must independently satisfy the instantiated
signature. This separation keeps the static rule simple, but shifts coverage
and implementation availability to runtime.

\subsection{Static environments and judgments}

We make the checker state explicit with three environments. Let \(\Gamma\)
map dynamic variables to types, let \(\Delta\) map template names to
signatures \((\vec p,T)\), and let \(\Xi\) map static variables in the
current scope to fresh type variables. We write
\[
  \Gamma;\Delta;\Xi \vdash e:T
  \qquad\text{and}\qquad
  \Xi \vdash_s T:T'
\]
for expression synthesis and static-type normalization, respectively. The
normalization judgment resolves bound names through \(\Xi\), treats other
names as constructors, normalizes \texttt{int} to \texttt{sint}, and checks
the arities of the built-in type constructors. Equality of normalized types
is structural: constructor names, function domains and codomains, and tuple
components must agree in order.

The literal, variable, abstraction, and application rules are conventional
for a monomorphic simply typed lambda calculus with annotated parameters:
\[
\frac{}{\Gamma;\Delta;\Xi\vdash n:\texttt{sint}}
\quad
\frac{\Gamma(x)=T}{\Gamma;\Delta;\Xi\vdash x:T}
\qquad
\frac{\Gamma,x:T_1;\Delta;\Xi\vdash e:T_2}
     {\Gamma;\Delta;\Xi\vdash\texttt{lam}(x:T_1)=>e:T_1\to T_2}
\]
\[
\frac{\Gamma;\Delta;\Xi\vdash e_1:T_1\to T_2
      \qquad \Gamma;\Delta;\Xi\vdash e_2:T_1}
     {\Gamma;\Delta;\Xi\vdash e_1(e_2):T_2}.
\]
The rules for Booleans, strings, tuples, projection, conditionals, and
operators enforce the constraints stated above: a conditional condition has
type \texttt{bool} and both branches have one common type; primitive
arithmetic takes \texttt{sint}; equality is limited to equal primitive
types; and projection selects an in-bounds component of a tuple. A recursive
function binds itself at the annotated type \(T_1\to T_2\), binds its
parameter at \(T_1\), and checks its body against \(T_2\).

For a declared signature \(\Delta(c)=(p_1,\ldots,p_n,T)\), static
application is checked by substitution, independently of the implementation
environment:
\[
\frac{\Delta(c)=(\vec p,T)
      \qquad \Xi\vdash_s T_i:T_i'\ (1\leq i\leq n)}
     {\Gamma;\Delta;\Xi\vdash \texttt{@}c[T_1,\ldots,T_n]:
       [T_1'/p_1,\ldots,T_n'/p_n]T}.
\]
The number of supplied arguments must equal the signature arity. A zero
parameter constant may also be written without brackets. Ordinary declarations
are sequential: a \texttt{val} initializer is checked before its name is
added to \(\Gamma\), and a \texttt{dec} extends \(\Delta\) only for later
declarations. A \texttt{let} checks its declarations in sequence and then
checks its body in the resulting environment; a bodyless \texttt{let} has
unit type.

For an implementation, let \(\Delta(c)=(\vec p,T)\), and let \(\vec q\)
be fresh rigid variables for its own parameters. The checker normalizes its
patterns \(\vec P\) under \(\Xi,\vec q\), requires every variable in
\(\vec q\) to occur in those patterns, and checks the implementation body
against the signature instantiated by those patterns:
\[
\frac{\Delta(c)=(\vec p,T)\quad
      \operatorname{vars}(\vec P)=\vec q\quad
      \Gamma;\Delta;\Xi,\vec q\vdash e:[\vec P/\vec p]T}
     {\Gamma;\Delta;\Xi\vdash
       \texttt{imp}\{\vec q\}\ c[\vec P]=e\ \mathsf{ok}}.
\]
Here \(\operatorname{vars}(\vec P)=\vec q\) abbreviates the requirement
that every declared implementation parameter occurs at least once; patterns
may of course contain repeated occurrences. The checker rejects duplicate
signature parameters, duplicate implementation parameters, duplicate visible
signatures, arity mismatches, unbound static variables, and bodies whose types
do not match the instantiated signature. It accepts overlapping patterns.

\subsection{Two environments}

Evaluation uses a value environment \(\rho\) and an implementation
environment \(\Theta\). The value environment maps variable names to values.
Function closures capture both \(\rho\) and \(\Theta\) at their definition
site. An implementation closure captures its definition-time value
environment, but template references appearing in its body use the active
implementation environment at the point of invocation. This distinction
preserves private value dependencies while allowing the active collection of
implementations to govern nested template calls.

Evaluation is eager and left-to-right, except that a conditional evaluates
only its chosen branch. Processing an implementation registers a closure and
does not execute its body. A template application instantiates and evaluates
the selected body afresh; the current interpreter does not memoize results.
The expression evaluator and type checker are separate entry points, so a
client is expected to check a program before evaluating it.

\subsection{Pattern matching and template resolution}

Static patterns are types in the same structural grammar: a variable, a
constructor application, a function type, or a tuple type. Write
\(P\matchsto T\theta\) when pattern \(P\) matches type \(T\) and produces
substitution \(\theta\). A variable binds on its first occurrence and must
be equal to its previous binding on every later occurrence. Constructors
match only constructors with the same name and arity; their arguments match
recursively. Function and tuple nodes match only the same node form, with
ordered children matched recursively. Matching the vector of patterns uses
one shared substitution:
\[
  (P_1,\ldots,P_n)\matchsto(T_1,\ldots,T_n)\theta.
\]
Consequently, \((a,a)\) matches \((\texttt{sint},\texttt{sint})\) with
\(\theta=\{a\mapsto\texttt{sint}\}\), but does not match
\((\texttt{sint},\texttt{bool})\). This is directional pattern matching,
not symmetric unification: variables in the left-hand pattern bind, while
right-hand variables are treated as literal type variables.

At runtime, \(\Theta=[K_1,\ldots,K_n]\) is a newest-first sequence of
implementation closures. Each closure \(K_i=(C_i,\rho_i)\) contains an
implementation declaration \(C_i\) and the value environment captured when
that implementation was registered. Write \(\mathit{name}(C_i)\),
\(\mathit{pat}(C_i)\), and \(\mathit{params}(C_i)\) for its constant name,
pattern vector, and declared implementation parameters.

Before resolution, requested types are normalized (including the
\texttt{int}/\texttt{sint} alias) and must be ground: no static variable may
remain. Define a candidate match by
\[
  \mathit{candidate}(K_i,c,\vec T,\theta)
  \quad\Longleftrightarrow\quad
  \mathit{name}(C_i)=c\ \land\
  \mathit{pat}(C_i)\matchsto\vec T\theta.
\]
The vector match requires equal lengths and threads one substitution through
all patterns. A pattern variable binds to the corresponding requested type
at its first occurrence. Each later occurrence succeeds only when the new
requested type is structurally equal to the existing binding. Constructor
nodes require equal names and arities; function and tuple nodes require the
same form and recursively matching ordered components. Matching is
directional: only variables in the left-hand pattern bind.

Resolution selects the least sequence index that is a candidate match:
\[
\frac{
  i=\min\{j\mid 1\leq j\leq n\ \land\
                 \exists\theta'.\,\mathit{candidate}(K_j,c,\vec T,\theta')\}
  \qquad
  \mathit{candidate}(K_i,c,\vec T,\theta)
  \qquad
  \operatorname{dom}(\theta)=\operatorname{set}(\mathit{params}(C_i))
}{\Theta\vdash c[\vec T]\Downarrow_R(K_i,\theta)}.
\]
Equivalently, every earlier closure either implements another constant or
fails to match the complete argument vector. Entries for other constants do
not affect priority. Well-typed programs satisfy the domain condition
because each implementation parameter must occur in a pattern. A malformed
directly constructed syntax tree that violates it raises a runtime error
rather than causing the evaluator to choose a later implementation.

If no closure matches, resolution has an explicit failure outcome:
\[
  \frac{\nexists i,\theta.\ \mathit{candidate}(K_i,c,\vec T,\theta)}
       {\Theta\vdash c[\vec T]\Downarrow_R\mathsf{no\ match}}.
\]
The interpreter reports this outcome as a missing-implementation runtime
error. Since the type checker does not inspect \(\Theta\), this failure is
possible even when the static application is well-typed.

After successful resolution, substitution \(\theta\) is applied to a fresh
copy of the selected body. Evaluation uses the selected closure's captured
value environment \(\rho_i\) and the active implementation environment
\(\Theta\):
\[
\frac{\Theta\vdash c[\vec T]\Downarrow_R((C_i,\rho_i),\theta)
      \qquad
      \langle\theta(C_i.\mathit{body}),\rho_i,\Theta\rangle\Downarrow v}
     {\langle\texttt{@}c[\vec T],\rho,\Theta\rangle\Downarrow v}.
\]
The caller's \(\rho\) does not replace the captured value environment.
Conversely, nested template references see active \(\Theta\), not a template
environment captured when \(C_i\) was registered. Bare \texttt{@c}
abbreviates \texttt{@c[]} for a zero-parameter constant.

Substitution traverses type annotations, nested static applications, and
declarations in the body. On entering a nested \texttt{dec} or \texttt{imp},
it removes mappings shadowed by that declaration's own static parameters.
Requested types are ground, so this shadowing rule prevents capture without
introducing free static variables. The source syntax tree is left unchanged;
each use gets its own instantiated copy.

Because \(\Theta\) is ordered newest first, dispatch priority follows source
registration order in reverse: a later matching implementation takes
precedence. The language currently has no overlap diagnostic or
specificity-based ordering.

\subsection{A resolution--evaluation commutation theorem}

For a successful resolution, write
\[
  \mathit{resolve}(\Theta,c,\vec T)=(K_i,\theta),
  \qquad K_i=(C_i,\rho_i),
\]
as an abbreviation for the successful resolution judgment
\(\Theta\vdash c[\vec T]\Downarrow_R(K_i,\theta)\),
and let \(\theta C_i.\mathit{body}\) denote capture-avoiding substitution
into a fresh copy of the selected implementation body. The order of
operations at a template reference is characterized by the following
local correspondence theorem. Here \(\Theta\) is the effective environment
at the occurrence being evaluated. It need not be the environment currently
held by the caller: the body of an ordinary function closure is evaluated
with the template environment captured when that function was created.

\medskip
\noindent\textbf{Theorem (template resolution commutes with evaluation).}
Suppose \(\vec T\) is ground, resolution is defined and binds exactly the
implementation parameters of \(C_i\), and both evaluations start with the
same active template environment \(\Theta\). Subjudgments may extend that
environment according to declaration scope. Then, for every value \(v\),
\[
  \langle\texttt{@}c[\vec T],\rho,\Theta\rangle\Downarrow v
  \quad\Longleftrightarrow\quad
  \langle\theta C_i.\mathit{body},\rho_i,\Theta\rangle\Downarrow v.
\]
In particular, resolving the reference before evaluating its selected body
preserves its result. The caller's value environment \(\rho\) is irrelevant
after resolution; the implementation's captured environment \(\rho_i\) is
used on both sides of the correspondence. \hfill\(\square\)

For example, if \texttt{f} is defined after \texttt{imp choose = 1}, its
closure captures a template environment containing that implementation.
Registering \texttt{imp choose = 2} later does not change the \(\Theta\)
used to evaluate \texttt{@choose} inside \texttt{f}. The theorem therefore
resolves that occurrence against the captured environment, where it selects
the first implementation. Resolution against the later top-level environment
would be a different judgment and is not covered by the theorem.

\noindent\emph{Proof sketch.} For the forward direction, invert the big-step
derivation for the template-reference expression. Its only successful rule
first resolves \(c[\vec T]\), substitutes the resulting bindings into the
selected body, and evaluates that body under \((\rho_i,\Theta)\). The
resolution relation is deterministic because it chooses the least matching
index and structural matching produces at most one binding for each pattern
variable. For the reverse direction, apply the template-reference evaluation
rule using the assumed resolution result and the given body evaluation. No
other successful resolution can be chosen. This theorem is a direct property
of the interpreter's semantics; it is not a proof of type preservation or
termination.\par

The result lifts to a demanded occurrence inside a call-by-value evaluation
context: replacing that occurrence by its selected, instantiated body leaves
the enclosing evaluation result unchanged, provided the surrounding context
and the active environments are held fixed. The demand condition matters.
The theorem does not authorize an eager whole-program pass that resolves
every syntactic reference before evaluation. A reference in an unselected
conditional branch is not evaluated, and a template reference inside a
function body is reached only when that function is called. Such a global
pass could report a missing implementation for code that the interpreter
never demands. It could also choose a different implementation if it failed
to preserve the environment at the occurrence.

\subsection{Big-step evaluation}

We write \(\langle e,\rho,\Theta\rangle\Downarrow v\) when expression
\(e\) evaluates to value \(v\). Values include base literals, tuples,
constructed values, and closures. The core closure rules are:
\[
\frac{}{\langle x,\rho,\Theta\rangle\Downarrow\rho(x)}
\qquad
\frac{}{\langle\texttt{lam}(x:T)=>e,\rho,\Theta\rangle
  \Downarrow \langle\texttt{lam}(x:T)=>e,\rho,\Theta\rangle}
\]
\[
\frac{\langle e_1,\rho,\Theta\rangle\Downarrow
          \langle\texttt{lam}(x:T)=>e,\rho_0,\Theta_0\rangle
      \quad \langle e_2,\rho,\Theta\rangle\Downarrow v_2
      \quad \langle e,\rho_0[x\mapsto v_2],\Theta_0\rangle\Downarrow v}
     {\langle e_1(e_2),\rho,\Theta\rangle\Downarrow v}.
\]
The recursive closure rule is analogous, with the closure itself bound to
the recursive name as well as the argument name. Since both names must be
distinct, the two bindings do not conflict. Calls evaluate the function and
then the argument before entering the closure; the body uses the closure's
captured environments.

The conditional rules evaluate the condition first and evaluate only the
selected branch. Tuple components and primitive operands evaluate left to
right. Integers use Python floor division and remainder behavior in the
prototype; division or modulo by zero is a runtime error. Runtime checks also
reject invalid operator operands, non-function application, invalid tuple
projections, and malformed constructed-value operations.

For declarations, a \texttt{val} evaluates its initializer and extends the
current value environment for later declarations. A \texttt{dec} has no
runtime effect. An \texttt{imp} prepends a closure containing the
implementation and a snapshot of the current value environment to
\(\Theta\); its body is not evaluated at registration time. Thus an
implementation can retain a private value from its definition site.
An implementation closure captures the current value environment when it is
registered; its body is not evaluated at registration time.

\section{Surface language}

The parser accepts a library as a sequence of declarations and a program as a
sequence of declarations followed by one result expression. Semicolons
separate declarations. The grammar below omits lexical details such as string
escapes and operator precedence, which are conventional and documented in
the implementation guide.

\begin{verbatim}
libfile       ::= declarations
program       ::= declarations expression
declarations  ::= [ declaration { ";" declaration } [";"] ]
declaration   ::= "val" id "=" expression
                | "dec" id ["[" [ids] "]"] ":" type
                | "imp" ["{" ids "}"] id ["[" [types] "]"]
                    "=" expression

expression    ::= "lam" "(" id ":" type ")" "=>" expression
                | "fix" id "(" id ":" type ")" ":" type
                    "=>" expression
                | "if" expression "then" expression "else" expression
                | "let" declarations ["in" expression] "end"
                | comparison
comparison    ::= additive [cmp additive]
additive      ::= multiplicative {("+" | "-") multiplicative}
multiplicative::= unary {("*" | "/" | "%") unary}
unary         ::= ("+" | "-") unary | postfix
postfix       ::= atom {"(" [expression-list] ")" | "." integer
                   | "[" [types] "]"}
atom          ::= integer | string | "true" | "false" | id | "@" id
                | generic-operator | parenthesized-expression
generic-operator ::= "D2CON" "(" string {"," expression} [","] ")"
                   | "D2TCK" "(" expression "," string ")"
                   | "D2ARG" "(" expression "," expression ")"

type          ::= type-atom ["->" type]
type-atom     ::= id [constructor-arguments] | parenthesized-type
constructor-arguments ::= "[" [types] "]"
                         | "(" [types] ")"
parenthesized-expression ::= "(" [expression
                              ["," [expression-list] ] ] ")"
parenthesized-type ::= "(" [type ["," [types]]] ")"
expression-list ::= expression {"," expression} [","]
types         ::= type {"," type} [","]
ids           ::= id {"," id} [","]
\end{verbatim}

In this compact grammar, \texttt{cmp} denotes one of \texttt{<},
\texttt{>}, \texttt{<=}, \texttt{>=}, \texttt{==}, or \texttt{!=};
parenthesized forms also distinguish grouping, unit, singleton tuples, and
tuples by the presence of commas. The parser gives calls, projections, and
static application a shared postfix precedence, then unary operators,
multiplication, addition, and non-chainable comparisons. A syntactic static
application may follow any postfix expression, but the checker and evaluator
currently require its target to be a template constant. The grammar describes
the accepted source shape; it does not imply that every syntactically
representable AST is well typed.

An identifier in a type is a static variable only when bound by an enclosing
signature or implementation parameter list; otherwise it is a named type
constructor. In expression position, an unprefixed identifier is a dynamic
variable, while \texttt{@id} denotes a template constant. A function accepts
one dynamic argument; a tuple packages multiple components into that single
argument. The parser builds an explicit level-2 syntax tree and does not
perform type checking or evaluation.

\section{Declaration checking}

The checking environment is extended from left to right. A value declaration
\texttt{val x = e} first synthesizes \(T\) for \(e\), then binds
\(x:T\). Thus a declaration cannot refer to its own newly introduced name,
although it may refer to earlier names. A signature
\texttt{dec c[\(\vec p\)]:T} creates fresh internal identities for its
parameters, normalizes \(T\) under those bindings, and records the resulting
signature. Fresh identities prevent accidental capture when nested scopes
reuse the same source spelling. A second visible declaration of the same
constant is rejected; dynamic variables may shadow earlier dynamic variables.

For \texttt{imp\{\(\vec q\)\} c[\(\vec P\)] = e}, the checker requires
that \(c\) already have a signature, that the pattern vector have the
signature's arity, and that every rigid implementation parameter occur in a
pattern. It substitutes the normalized patterns for the signature
parameters, then checks \(e\) against the resulting type under the
implementation's static scope. This is an independent check of each
implementation; it neither checks dispatch coverage nor rejects overlaps.

Declarations are processed sequentially in one environment. The checker has
no declaration form for introducing a nested visibility region.

\section{Worked examples}

\subsection{Generic factorial over a template interface}

The numeric library declares constants for zero and one, conversion from an
integer, comparisons, and arithmetic. The core factorial implementation is:
\begin{lstlisting}
dec f0_1[a]: () -> a;
dec f1_lte[a]: (a, a) -> bool;
dec f1_sub[a]: (a, a) -> a;
dec f1_mul[a]: (a, a) -> a;

dec f1_fact[a]: a -> a;
imp{a} f1_fact[a] =
let
  val one = @f0_1[a]()
in
  fix fact(n: a): a =>
    if @f1_lte[a](n, one) then one
    else @f1_mul[a](n, fact(@f1_sub[a](n, one)))
end;
\end{lstlisting}
The type checker verifies the body using only the signatures of
\texttt{f0\_1}, \texttt{f1\_lte}, \texttt{f1\_sub}, and
\texttt{f1\_mul}. The implementation of factorial is consequently valid for
any static type \(a\) for which those constants are declared, without
assuming that \(a\) is an integer. To run it, matching numeric
implementations must exist for the requested type. For \texttt{sint},
\texttt{@f1\_fact[sint](5)} evaluates to \texttt{120}. This example
illustrates the distinction between signature-level checking and run-time
availability of a whole family of operations.

The types do not establish the algebraic laws expected of the numeric
operations, nor that the input is nonnegative. The supplied recursive
definition does not memoize results, and its termination is not proved by the
checker.

\subsection{Overlapping patterns and priority}

Consider two implementations of the same signature:
\begin{lstlisting}
dec choose[a,b]: () -> sint;
imp{x,y} choose[x,y] = lam(u: ()) => 1;
imp{x} choose[x,x] = lam(u: ()) => 2;
\end{lstlisting}
Both bodies satisfy the signature. The second implementation is registered
later and is therefore examined first. It matches equal arguments, so
\texttt{@choose[sint,sint]()} returns \texttt{2}; for unequal arguments its
pattern fails and the earlier general implementation returns \texttt{1}.
Reversing the declarations changes the result for equal arguments. This is a
deterministic ordered-dispatch policy, but it is sensitive to declaration
order and is not a coherence or most-specific-instance discipline.

\subsection{Algebraic values and streams}

The libraries represent lists and options as named constructor values. The
generic expression \texttt{D2CON("list\_cons", h, t)} creates a value with an
ordered pair of fields; \texttt{D2TCK(v,"list\_cons")} tests its tag, and
\texttt{D2ARG(v,0)} extracts its head. The static checker recognizes the
supported list, option, and stream constructor cases and checks their
arguments. These operations do not introduce a user-extensible datatype
declaration mechanism, and the interpreter represents constructor identity
with runtime names and ordered fields.

A stream of elements of type \(a\) is represented as a delayed unit function
that returns a stream-constructor node. Applying a stream to \texttt{()} is
the force operation; its result has type \texttt{stcn(a)}. The supplied
stream library builds the delayed node with \texttt{strm\_cons}. Forcing is
not memoized: evaluating the thunk again recomputes the node. This is a
concrete representation convention enforced by the checker for stream
application, rather than a general lazy evaluation strategy.

\section{Implementation and validation status}

The prototype is written in Python 3.12 using only the standard library. Its
front end parses concrete source into level-2 dataclasses. Static types and
expressions share a small syntax-tree representation; source locations are
retained by parser errors but are not stored on the resulting trees. The
checker produces an independent type environment and does not rewrite the
program. The interpreter then evaluates the same tree using explicit closure
records and immutable snapshots of the variable and template environments.

The implementation is divided into modules for AST definitions, parsing,
structural equality and matching, type checking, and interpretation. Library
files demonstrate integer operations, generic factorial and Fibonacci, list
and option templates, and stream constructors. The repository includes
separate unit-test suites for parsing, type checking, and interpretation. The
test descriptions cover static scope, generic implementations, nested
substitution, private closure dependencies, overlapping-pattern order, and
example-library execution. This draft has not measured performance or
compared the interpreter against a compiler; no such claims are made.

The execution pipeline is therefore:
\[
  \text{source}\longrightarrow\text{AST}
  \longrightarrow\text{type check}
  \longrightarrow\text{interpret}.
\]
The stages are callable independently. In particular, the interpreter does
not invoke the checker automatically. A client that evaluates unchecked
syntax may encounter runtime errors for unbound variables, absent template
implementations, malformed constructor access, invalid arithmetic operands,
or division by zero.

\section{Relation to prior work}

\FWATS{} uses terminology and design concerns that also occur in ATS,
including static parameters, templates, and the distinction between static
and dynamic computation. ATS is a substantially richer language and type
system: its formal foundations treat statics and dynamics as separate
components and study dependent types and programming with theorem-proving
\cite{xi-ats-foundation}. The ATS programming materials also describe
template-based genericity and late-binding in a compiler-oriented setting
\cite{ats-template-doc}. The present prototype should not be read as an ATS
implementation or a formal subset of ATS. Its static language is limited to
structural type expressions, and its templates are resolved by a small
run-time interpreter over explicitly supplied ground arguments.

At a broad level, the signature/implementation split also invites comparison
with type classes and modular dispatch systems. The operational distinction
here is explicit and intentionally narrow: the checker uses a declared
signature without selecting an implementation, while the interpreter
selects the newest matching implementation in a runtime environment. This
draft does not claim that the mechanism is novel relative to those systems;
establishing a more exact comparison requires a focused study of coherence,
instance selection, and separate compilation in the relevant languages.

\section{Limitations and next steps}

Several omissions define the current research boundary. Static arguments are
explicit; template values are not first-class; constructor declarations are
built in rather than user-defined; and no inference or specificity ordering
is provided. The checker does not guarantee a matching implementation,
termination, absence of division by zero, or semantic laws for user-supplied
operations. Integer division follows Python floor-division semantics.
Recursive evaluation is subject to Python's recursion limit and has no tail
call optimization. Source locations are not available to type and runtime
errors after parsing.

A fuller formal account should next establish metatheoretic properties for a
precisely delimited core, or explicitly characterize which failures are
intended runtime outcomes. In particular, useful extensions include a
well-formedness condition connecting checked declarations to runtime
implementation environments, a proof that successful resolution preserves
the declared result type, diagnostics for overlapping patterns, and a
mechanized account of closure capture. Any change from
ordered dispatch to coherence or specificity-based selection would alter
observable program behavior and should be treated as a language-design
decision.

\section{Conclusion}

This draft has described FWATS3 as a typed expression language with explicit
static arguments and separately registered implementations. Its operational
model makes template resolution a first-match search over an ordered runtime
environment, while its type checker validates implementation bodies against
instantiated signatures without proving implementation coverage. The result
is a compact executable setting in which to study the interaction of
structural type checking, pattern-directed specialization, lexical value
capture, and dynamic template resolution. The current implementation is a
prototype rather than a verified language implementation; the semantics and
limitations stated here are intended to guide a more complete formalization.

\begin{thebibliography}{9}
\bibitem{xi-ats-foundation}
Hongwei Xi.
\newblock Applied Type System: An Approach to Practical Programming with
Theorem-Proving.
\newblock Manuscript on the formal foundations of ATS, 2016.
\newblock \url{https://www.cs.bu.edu/~hwxi/atslangweb/MYDATA/ATSfoundation.pdf}.

\bibitem{ats-template-doc}
ATS Documentation Project.
\newblock Introduction to Programming in ATS, section on parametric
polymorphism and template-based programming.
\newblock \url{https://ats-lang.github.io/FROZEN000/DOCUMENT/INT2PROGINATS/HTML/HTMLTOC/c1126.html}.
\end{thebibliography}

\end{document}
